% Issue #81, part "Architecture logic"; owner: 11yashiagrawal. Input from
% reports/phase2/main.tex (tikz with arrows.meta and positioning is loaded there).
% Numbers: results/metrics/fd001_{lstm,tcn,transformer}_grid.csv (grids, weights)
% and fd001_cv_{lstm,tcn,transformer}{,_uncapped}.json (chosen models).
% Settings: configs/{dl_base,lstm,tcn,transformer}.yaml; contract: src/rul/models/base.py.

\section{Model Architectures}
\label{sec:architecture}

Phase~2 asks whether reading the raw sensor sequence beats the Phase~1
rolling features, and, if it does, which way of reading it works best. To
answer the second question fairly, the three encoders differ in one respect
only: how information moves across the window. Everything else is shared.

\subsection{The Shared Contract}
Every sequence model implements one interface
(\texttt{rul.models.base}, Fig.~\ref{fig:contract}). An \emph{encoder} maps
a batch of windows of shape $[B,L,F]$ ($L=30$ cycles, $F$ scaled sensors
chosen on the training fold) to one embedding per window, $[B,d]$. A shared
regression \emph{head} (a linear layer) maps the embedding to one RUL value.
The contract fixes the rest of the pipeline as well. All models use the same
windows and padding (Section~\ref{sec:dl_dataset}), the same capped target and
selection metric (Section~\ref{sec:target_selection}), the same trainer (MSE,
Adam, cosine schedule, gradient clipping, early stopping on held-out engines,
seed 42), the same three engine-grouped folds as Phase~1, and the same metric
code.

This design has three consequences. First, a difference in score can be
attributed to the encoder, because nothing else changes. Second, ablations
can swap a single component. Third, the hybrids reuse any encoder as their
sequence branch without new training code (Section~\ref{sec:theory}).

\begin{figure}[t]
\centering
\begin{tikzpicture}[
  box/.style={draw, rounded corners=2pt, align=center, font=\footnotesize,
              minimum height=8mm, inner sep=3pt},
  lab/.style={font=\scriptsize, midway, above},
  >={Stealth[length=2mm]}]
\node[box] (x) {window\\$[B,30,F]$};
\node[box, right=7mm of x] (enc) {encoder\\LSTM / TCN /\\Transformer};
\node[box, right=7mm of enc] (head) {linear\\head};
\node[box, right=6mm of head] (y) {RUL\\$[B]$};
\draw[->] (x) -- (enc);
\draw[->] (enc) -- node[lab]{$[B,d]$} (head);
\draw[->] (head) -- (y);
\node[font=\scriptsize, align=center, below=3mm of enc]
  {shared: windows, cap, folds, trainer, seed, metrics};
\end{tikzpicture}
\caption{The contract every Phase~2 sequence model implements. Only the
encoder changes between the LSTM, the temporal CNN and the Transformer.}
\label{fig:contract}
\end{figure}

\subsection{Why These Three Encoders}
Each encoder makes a different bet about how wear shows up in a 30-cycle
window. Section~\ref{sec:theory} gives the mechanics.

\emph{LSTM: wear is a slow trend.} The recurrent cell reads the cycles in
order and keeps a gated memory, so a drift spread over many cycles can build
up while single-cycle noise is filtered out. It is the standard model on
C-MAPSS~\cite{Zheng2017long,Chen2021machine}, which makes it the reference
point.

\emph{Temporal CNN: wear is a local pattern.} Dilated causal convolutions
look for short shapes, such as a step or a change of slope, anywhere in the
window, and process all positions in parallel. Its receptive field is a known
number rather than a learned memory: our setting sees 61 cycles, enough for
the whole window.

\emph{Transformer: wear is a relation between distant cycles.}
Self-attention compares every cycle with every other cycle directly, for
example the latest reading against the start of the window, without passing
information through intermediate steps. It assumes the least about structure,
which is also why it must stay small on 100 training engines.

Together the three cover the main options: recurrence, convolution and
attention. A result such as ``the LSTM wins'' is therefore a statement about
which assumption suits FD001, not about one tuned network.

\subsection{How Each Was Configured}
Each encoder had its own grid of 12 to 16 settings, scored with the
selection RMSE (held-out cycles with true RUL $\leq130$) over the three folds.
The LSTM keeps its best setting. For the temporal CNN and the Transformer we
took the \emph{smallest} setting within one fold standard deviation of the
best. A difference smaller than the fold-to-fold spread is not evidence, and
with only 100 engines the smaller model is the safer bet.
Table~\ref{tab:arch_choice} lists the chosen settings.

\begin{table}[t]
\centering
\caption{Chosen encoders, FD001 cross-validation (mean $\pm$ s.d.\ over 3 folds,
seed 42). Weights include the head; the Transformer has one layer. ``$\leq130$'': selection RMSE on held-out
cycles with true RUL $\leq130$, for capped and uncapped training.}
\label{tab:arch_choice}
\footnotesize
\setlength{\tabcolsep}{3pt}
\begin{tabular}{lrcc}
\hline
Encoder & Weights & Capped $\leq130$ & Uncapped $\leq130$ \\
\hline
LSTM, $1\times64$ & 20,801 & $15.78\pm0.44$ & $30.79\pm1.51$ \\
TCN, 32 ch., $k{=}3$, 4 lev. & 23,745 & $17.24\pm1.19$ & $29.93\pm2.47$ \\
Transf., $d{=}32$, 4 heads & 9,153 & $15.59\pm0.23$ & $30.99\pm3.23$ \\
\hline
\end{tabular}
\end{table}

Three points stand out.
\begin{enumerate}
\item \emph{The target matters more than the architecture.} Under the cap,
the three encoders are within 1.7 cycles of each other. Trained on the
uncapped target, all three land near 30. Removing the cap costs each encoder
13 to 15 cycles, against 1.7 cycles between the best and worst encoder, so on
FD001 the training target matters far more than the architecture.
\item \emph{More capacity did not help.} Every two-layer LSTM scored worse
than the chosen one-layer model, and so did 128 units against 64. The best
TCN (149,761 weights) and the best Transformer (34,689) were each within one
fold standard deviation of the much smaller settings we chose. With 100
engines, data, not capacity, limits these models.
\item \emph{The TCN is the least stable.} Its fold spread (1.19) is the
largest of the three, and its fitting was the slowest (306 s, against 57 s for
the LSTM and 117 s for the Transformer on CPU).
\end{enumerate}
These are cross-validation scores used for choosing models. The comparison on
the official test set comes after the model freeze (\#71) and is reported
with the results.
